\documentclass[11pt]{article}
\usepackage[margin=1in]{geometry}
\usepackage{float}
\usepackage{amsmath,amssymb,bm,booktabs,graphicx,hyperref}
\hypersetup{colorlinks=true,linkcolor=blue,urlcolor=blue}

\title{CS 577 Deep Learning\\Homework 1: Tools and Tensor Fluency}
\author{Your Name (Hawk username: \texttt{Irma Modzgvrishvili})}
\date{Fall 2026}

\begin{document}
\maketitle

\section*{Part 1: LaTeX and Mathematical Communication}

\subsection*{1.1 Tensor notation}

Typest the following statement cleanly, including dimensions:
\begin{quote}
For a batch $X$ with shape $(B,D)$, weights $W$ with shape $(D,K)$,
and bias $b$ with shape $(K,)$, the affine output is

\[
Y = XW + 1_Bb^T
\]

and has shape $(B,K)$
\end{quote}
\noindent
Here, $1_B \in \mathbb{R}^{B \times 1}$ denotes the all-ones column vector. Therefore, $1_Bb^T \in \mathbb{R}^{B \times K}$ repeats the bias vector $b^T$ once for each of the $B$ examples. This is the mathematical form of broadcasting $b$ across the batch dimension in the NumPy/PyTorch expression 

\[
Y = XW + b
\]
\noindent
Use mathematical notation rather than copying the monospace statement verbatim.

\subsection*{1.2 Regularized mean-squared error}
Typest the objective
\[
J(w)=\frac{1}{N}\|Xw-y\|_2^2+\frac{\lambda}{2}\|w\|_2^2
\]
\noindent
and its gradient with respect to $w$

\[
  \nabla_w J(w) = \frac{2}{N} X^T (Xw - y) + \lambda w
\]
\subsection*{1.3 Results table and interpretation}
\begin{table}[H]
  \centering
  \begin{tabular}{lr}
    \toprule
    Quantity & Value \\
    \midrule
    Initial training loss & 2.347500 \\
    Final training loss & 1.257789e-06 \\
    Learned slope & 2.498237 \\
    Learned intercept & -0.399999 \\
    Maximum gradient difference & 3.100408818568212e-12 \\
    \bottomrule
  \end{tabular}
  \caption{Results from the deterministic regression and gradient checks.}
\end{table}

% Add a 3-5 sentence interpretation here.
\noindent\textbf{Interpretation:}
The model optimization run successfully because training loss drop from initial value 2.3475 down to final loss 1.25779. The algorithm learn a slope around 2.49823 and intercept around -0.39999 to fit target dataset. Also, the maximum gradient difference between numeric check and math formula is extremely small. This confirm that our analytical gradient derivation is implemented correct.

\section*{Part 2: NumPy Tensors and Vectorization}

\subsection*{2.1 Shapes and broadcasting}
% State shapes, explain broadcasting, and report Y.
shapes: X.shape=$(4, 3)$
W.shape= $(3, 2)$ 
b.shape=$(2,)$ 
Y.shape=$(4, 2)$
\\
$X @ W$ is matrix multiplication. Since $X$ has shape $(4,3)$ and $W$ has shape $(3,2)$, the inner dimensions match $(3 = 3)$, so the result has shape $(4,2)$. The bias $b$ has shape $(2,)$. Broadcasting compares dimensions from right to left, so $(2,)$ is treated like $(1,2)$. Since $2 = 2$ and 1 can expand to 4, $b$ is broadcast to shape $(4,2)$ and added to every row of $X @ W$.
\[
Y =
\begin{bmatrix}
4.25 & -3.75 \\
0.25 & 5.25 \\
-2.25 & -0.75 \\
-1.25 & 4.25
\end{bmatrix}
\]


\subsection*{2.2 Standardization}
% Report mean, standard deviation, and checks on Z.

\noindent\textbf{mean:}
\[
[0.5,\ 0.25,\ 1.25]
\]

\noindent\textbf{std:}
\[
[1.11803399,\ 1.47901995,\ 1.47901995]
\]

\noindent\textbf{check mean:}
\[
[0.0,\ -1.38777878\times10^{-17},\ 0.0]
\]

\noindent\textbf{check std:}
\[
[1,\ 1,\ 1]
\]


\subsection*{2.3 Pairwise squared distances}
% Report D2 and explain the inserted singleton dimensions.

\[
D2=
\begin{bmatrix}
2 & 1\\
1 & 4 \\
2 & 5
\end{bmatrix}
\]

\noindent
The inserted singleton dimensions allow broadcasting to compare every row of A with every row of B simultaneously, without a Python loop.
\subsection*{2.4 Vectorization timing}
% Report the benchmark setup, median times, and a careful interpretation.
We test matrix multiplication using explicit Python loops vs vectorized NumPy code on $1000 \times 1000$ matrices. We run 10 iterations and report median time:Python loops time: 0.0948s secondsVectorized NumPy time: 0.0019s secondsSpeedup factor: $\approx 49.86x\times$ fasterInterpretation:Vectorized NumPy code is much faster cuz it use optimized C libraries and parallel CPU operations under the hood. Using explicit Python for loops is very slow for large matrix because of Python execution overhead.
\section*{Part 3: PyTorch Tensors and a Minimal Model}

\subsection*{3.1--3.2 Agreement, dtype, and device}
Tensors participating in an operation generally need to be on the same device because CPU and GPU memory are separate. Their dtypes also need to be compatible; otherwise PyTorch may raise an error or perform a type conversion. Explicitly choosing both the device and dtype makes the computation predictable
\subsection*{3.3--3.4 Model and training}
% Summarize architecture and training results; detailed code belongs in the notebook.
We build a simple 2-layer network using PyTorch (Linear $\rightarrow$ ReLU $\rightarrow$ Linear). 

\begin{itemize}
    \item \textbf{Architecture:} Input dim = 10, Hidden dim = 32, Output dim = 1.
    \item \textbf{Loss \& Optimizer:} Mean Squared Error (MSE) loss with Adam optimizer (learning rate = 0.01).
    \item \textbf{Training results:} Model trained for 100 epochs. Initial loss started at 1.452 and smoothly dropped down to 0.012, showing the model learned target mapping successfully.
\end{itemize}

\section*{Part 4: Three Views of the Same Gradient}

\subsection*{4.1 Analytic gradient}
% Derive the gradient.
For loss function
  \[
    f(w) = \frac{1}{2N} \Vert{}Xw - y\Vert{}_2^2
  \]
 we take derivative with respect to $w$:$$\nabla_w f(w) = \frac{1}{N} X^T (Xw - y)$$We use this formula in code to get exact math gradient vector.

\subsection*{4.2--4.3 Numerical comparison}
% Report the three vectors, errors, and explanation.
We calculate loss function and gradient at point $w$ using three different way. 
\begin{table}[h]
   \centering
   \begin{tabular}{ll}
    \toprule
    Metric & Value / Result \\
    \midrule
    Objective value & 0.2031666666666667 \\
    Analytic gradient & [-0.88, -0.16333333] \\
    Finite difference gradient & [-0.88, -0.16333333] \\
    Autograd gradient & [-0.88, -0.16333333] \\
    \midrule
    \multicolumn{2}{l}{\textbf{Max Absolute Pairwise Differences:}} \\
    Analytic vs Finite difference & 3.100408818568212e-12 \\
    Analytic vs Autograd & 0.0 \\
    Finite difference vs Autograd & 3.100408818568212e-12 \\
    \bottomrule
   \end{tabular}
\end{table}
The analytic gradient and PyTorch autograd match perfectly (0.0 difference) cuz autograd calculate exact derivatives using chain rule. The finite difference method has a super small error ($\sim 10^{-12}$) cuz it is a numerical approximation using small step size $\epsilon$. This tiny difference proves our math code and autograd implementation are both correct.

\section*{Part 5: Reflection}

\paragraph{Shape/broadcasting mistake.} When doing squared distance, forgetting to use A[:, None, :] and B[None, :, :] break the shapes and broadcasting doesnt work right.

\paragraph{Gradient validation method.} Using finite difference formula 
\[
\frac{f(w + \epsilon) - f(w - \epsilon)}{2\epsilon}
\]
 gives us a easy way to check if the math gradient and PyTorch autograd are correct.

\paragraph{Workflow change before HW2.} Before HW2, I will always run unittest script first to catch small bugs early.

\section*{Assistance and Generative-AI Disclosure}

% Name each tool/person, purpose, and what you verified or changed. If none, state that.
TODO

\end{document}

